\documentclass{article}
\usepackage[T1]{fontenc}
\usepackage{lmodern}
\usepackage{graphicx}
\usepackage{amsmath,amssymb}
\usepackage[a4paper,margin=2cm]{geometry}
\usepackage[absolute,overlay]{textpos}
\setlength{\TPHorizModule}{1mm}
\setlength{\TPVertModule}{1mm}
\usepackage[indonesian]{babel}
\usepackage{hyperref}
\setlength{\emergencystretch}{2em}
% unit: Halaman 57

\begin{document}

% === Halaman 57 (python/native) ===

57

Ekstraksi pesan dari Stego-image

• Posisi pixel yang menyimpan bit pesan dapat diketahui dari 

bilangan acak yang dibangkitkan oleh PRNG. 

• Jika kunci yang digunakan pada waktu ekstraksi sama dengan kunci 

pada waktu penyisipan, maka bilangan acak yang dibangkitkan juga 

sama. 

• Dengan demikian, bit-bit pesan yang bertaburan di dalam citra 

dapat dikumpulkan kembali.

\end{document}
